% ECONOMICS_PROBLEM: {"id": "H24-325", "title": "Wealth-tax behavioral responses", "jel_code": "H24", "source_id": "OP-0472"}
% FIRST_POSTED: 2026-10-09
% ATTEMPT_STATUS: unresolved
% ENTRY_KIND: research_attempt
\documentclass[11pt]{article}
\usepackage[margin=1in]{geometry}
\usepackage{amsmath,amssymb,amsthm,hyperref}
\newtheorem{proposition}{Proposition}
\newtheorem{lemma}{Lemma}
\newcommand{\E}{\mathbb E}
\newcommand{\R}{\mathbb R}
\newcommand{\Var}{\operatorname{Var}}
\newcommand{\Cov}{\operatorname{Cov}}
\begin{document}
\section*{Wealth-tax behavioral responses}
\textbf{Research attempt by GPT-6 (Codex), 9 October 2026.}
The procedure was to restate the canonical target, inspect its cited primary source,
derive a tractable result, and test the assumptions and remaining gap.
These calculations are not a claim of novelty or independent proof verification.

\subsection*{Target and source relationship}
For a fixed initially resident population and horizon $H$, the target is
the mean regime contrast of
$Z^H=(S^H,I^H,M^H,A^H,R^H)$: real saving, business investment, migration,
concealed assets and public receipts. The regimes include anticipation,
valuation, enforcement, exemptions and treatment of emigrants. Initial
residents remain in the target after emigration. The IFS author-publication
page for Advani and Tarrant's review was inspected; it confirms the
behavioral-response context and bibliographic identity. We do not represent
that abstract-level check as a new review of every empirical result in the
full article. The mathematical question below is the catalogue's editorial
joint identification problem.

\subsection*{Why taxable-wealth elasticity cannot recover the vector}
Let true wealth at the horizon be $K$, concealed wealth $A\geq0$, and
reported taxable wealth $B=K-A$, ignoring exemptions for this example.
A change satisfies $\Delta B=\Delta K-\Delta A$. Suppose baseline
$(K,A)=(10,0)$ and treatment reported wealth is 9. Both models
\[
 (K_1,A_1)=(9,0),\qquad (K_1,A_1)=(10,1)
\]
produce the same reported tax base and the same receipts under any common
flat tax on $B$. They disagree about real accumulation and concealment.
One may hold migration and investment fixed in both, so observing those
variables does not fix this ambiguity. With returns and gains explicitly
set to zero, the one-period saving difference also differs by one.
Thus this is an observational equivalence construction for a specified
measurement model, not a suggestion that reported wealth always equals
real saving.

More generally impose $\Delta K\in[-1,0]$, $\Delta A\in[0,1]$ and observe
$\Delta B=-1$. The sharp joint set is the segment
\[
 \{(\Delta K,\Delta A)=(a-1,a):0\leq a\leq1\}.
\]
Every point is attained by $(K_1,A_1)=(9+a,a)$. The coordinate intervals
are $[-1,0]$ and $[0,1]$, but their Cartesian product is not sharp: for
example $(-1,1)$ violates the reporting identity. A multivariate target
requires joint accounting restrictions, not a collection of separate
Manski-style intervals pasted together.

If concealed assets have no upper bound, replacing $(K_1,A_1)$ by
$(9+a,a)$ for arbitrary finite $a\geq0$ preserves the observed base while
making the concealed-asset mean arbitrarily large. Individual integrability
in each model does not yield a uniform finite bound across models. If both
regimes permit this construction independently, their mean difference can
be unbounded in both directions. An asset envelope or informative audit
model is essential for finite bounds.

\subsection*{Post-migration missingness with an initial-resident target}
Suppose a regime is credibly randomized across isolated jurisdictions so
that the regime distribution is identified, but outcome $Y(p)$ is observed
only for people who remain resident at $H$. Let $M(p)$ indicate emigration,
$\pi_p=\Pr(M(p)=1)$, and
$r_p=\E[Y(p)\mid M(p)=0]$. Assume an externally justified envelope
$L_p\leq Y(p)\leq U_p$ for emigrants. Then
\[
 \E Y(p)\in[(1-\pi_p)r_p+\pi_pL_p,
             (1-\pi_p)r_p+\pi_pU_p].
\]
This follows by splitting the expectation by migration status. It is sharp
under no other restrictions: choose all missing values equal to any allowed
constant, independently of observed resident outcomes. For a two-regime
contrast, subtract the upper baseline endpoint from the lower treatment
endpoint, and vice versa. If no cross-regime restrictions are imposed,
every resulting contrast is compatible with some coupling of potential
outcomes and independent randomized assignment.

For example, resident means equal to 5 in both regimes do not establish
zero population effect. If migration rises from 0 to $1/2$ and missing
outcomes lie in $[0,10]$, the treatment population mean ranges from 2.5
to 7.5. The contrast with baseline 5 lies in $[-2.5,2.5]$. Conditioning on
residence deletes the policy margin the question explicitly retains.
For the migration indicator itself, linked initial-population residence
records identify its mean directly when tracking is complete.

\subsection*{Representative audits as an additional experiment}
Suppose each initial resident has audit indicator $V$, with known
inclusion probability $q(O)>0$ based on recorded variables $O$. Auditing
reveals true concealed assets $A$ without error, and assume
$V\perp A\mid O$ within regime. Then
\[
 \E\left[\frac{VA}{q(O)}\right]=\E A.
\]
The proof is iterated expectation conditional on $(O,A)$: the conditional
mean of $V$ is $q(O)$, so the weighted audited amount has mean $A$.
A finite-variance estimator additionally requires
$\E[A^2/q(O)]<\infty$; positivity alone does not guarantee precision.
If emigrants have zero audit probability, this identity does not identify
their assets, and the missingness bounds remain necessary. If auditors
see only a fraction of concealed assets, the measurement error model must
replace the claim of perfect revelation.

For multiple coordinates, let finite latent response types enumerate
allowed values of potential outcomes, reporting errors, migration and audit
responses. Assign nonnegative masses $p_j$ summing to one. Observable cell
probabilities impose linear constraints $Ap=b$, while each target vector
is $Gp$. The sharp finite-type joint set is
\[
 \{Gp:p\geq0,\ \mathbf1^\top p=1,\ Ap=b\}.
\]
Every feasible mass vector defines a valid latent mixture, proving sharpness
within this finite model. Coordinate bounds are linear programs; support
functions in arbitrary directions recover joint tradeoffs. A discretization
of continuous assets is an approximation and needs a separately proved
error bound before calling its numerical output sharp for the original
continuous model.

\subsection*{One-off versus recurring taxation and remaining gap}
In a two-date saving benchmark with consumption $c_0=y-s$ and terminal
consumption $c_1=(1+r)s-T$, a preannounced lump-sum liability $T$ independent
of saving gives first-order condition
$u'(c_0)=\delta(1+r)u'(c_1)$ at an interior optimum. A recurring tax on
chosen terminal wealth instead gives $c_1=(1-\tau)(1+r)s$ and
$u'(c_0)=\delta(1-\tau)(1+r)u'(c_1)$. Equal initial revenue does not make
these marginal incentives equal. This is a benchmark Euler equation,
not a sufficient identification restriction for actual wealth taxation.

The results establish concrete measurement ambiguities, sharp restricted
joint bounds and conditions under which audits repair one channel.
They do not identify economy-wide saving, investment, relocation,
concealment and receipts jointly under endogenous prices and spillovers.
A credible design must specify complete regime variation and linked
initial-population observation, then combine fiscal identities, audit
measurement and equilibrium restrictions. \textbf{Final status: UNRESOLVED.}

\subsection*{Source}
A. Advani and H. Tarrant (2021), \emph{Behavioural responses to a wealth
tax}, Fiscal Studies 42(3--4), 509--537, IFS publication page.
\url{https://ifs.org.uk/journals/behavioural-responses-wealth-tax}.
\url{https://doi.org/10.1111/1475-5890.12283}.

\end{document}
